% Part: sets-functions-relations
% Chapter: arithmetization
% Section: checking-details
%
\documentclass[../../../include/open-logic-section]{subfiles}


\begin{document}

\olfileid{sfr}{arith}{check}
\olsection{Ring lan Lapangan Katata}

Sajrone bab iki, kita ngaku yen definisi-definisi tartamtu tumindak
``kaya kang kudune''. Ing lampiran teknis iki, kita bakal njlentrehake
apa kang dimaksud lan menehi gambaran bukti yen definisi-definisi iku
pancen tumindak ``kanthi bener''.

Ing \olref[int]{sec}, kita netepake panambahan lan ping-pingan ing $\Int$.
Kita kepengin nuduhake yen operasi kang wis ditetepake iki menehi struktur
kang ``dikarepake'' marang $\Int$. Mligine, struktur kasebut yaiku ring
komutatif.

\begin{defn}
	\emph{Ring komutatif} yaiku himpunan $S$ kang dilengkapi unsur tartamtu
	$0$ lan $1$ sarta operasi $+$ lan $\times$, kang nyukupi wolung rumus iki:
	\begin{align*}
		\emph{Asosiativitas}&&a + (b+ c) & = (a + b) + c \\
		&& (a \times b) \times c & = a \times (b\times c)\\
		\emph{Komutativitas}&&a + b &= b+ a  \\
		&&  a \times b&= b\times a\\
		\emph{Identitas}&&a + 0 &= a \\
		&& a \times 1 &= a\\
		\emph{Invers Aditif}&&(\exists b\in S)0&=a + b\\
		\emph{Distributivitas}&&a \times (b+ c ) &= (a \times b) + (a \times c)
	\end{align*}
	Kanthi implisit, kabeh variabel ing rumus iki kaiket dening kuantor
	universal kang diwatesi ing $S$. Gatekna uga yen unsur $0$ lan~$1$ ing
	kene ora kudu dadi wilangan natural kang jenenge padha.
\end{defn}

Dadi, kanggo mriksa yen wilangan wutuh mbentuk ring komutatif, kita mung
perlu mriksa yen wolung sarat iki kacukupan. Ora ana sarat kang {angel}
kanggo dibuktekake, nanging pamriksane rada ngrepotake. Minangka tuladha,
iki cara mbuktekake \emph{Asosiativitas} kanggo panambahan:

\begin{proof}
Tetepna $i, j, k \in \Int$. Dadi ana $a_1, b_1, a_2, b_2, a_3, b_3 \in
\Nat$ saengga $i = \equivrep{a_1, b_1}{}$ lan $j =
\equivrep{a_2,b_2}{}$ lan $k = \equivrep{a_3, b_3}{}$. (Supaya gampang
diwaca, kita nulis ``$\equivrep{x, y}{}$'' minangka ganti
``$\equivrep{x, y}{\Intequiv}$''; iki bakal ditindakake ing saindhenging
perangan iki.) Saiki:
\begin{align*}
	i + (j + k) &= \equivrep{a_1, b_1}{}+(\equivrep{a_2, b_2}{} + \equivrep{a_3, b_3}{}) \\
	&= \equivrep{a_1,  b_1}{} + \equivrep{a_2+a_3, b_2+b_3}{}\\
	&= \equivrep{a_1 + (a_2 + a_3), b_1 + (b_2 + b_3)}{}\\
	&= \equivrep{(a_1 + a_2) + a_3, (b_1 + b_2) + b_3}{}\\
	&= \equivrep{a_1 + a_2, b_1 + b_2}{} + \equivrep{a_3, b_3}{}\\
	&= (\equivrep{a_1, b_1}{} + \equivrep{a_2, b_2}{}) + \equivrep{a_3, b_3}{}\\
	&= (i+j) + k
\end{align*}
kanthi nggunakake tumindake panambahan ing $\Nat$ kanthi bebas.
\end{proof}

Kanthi cara kang padha, iki bukti \emph{Invers Aditif}:

\begin{proof}
Tetepna $i \in \Int$, mula $i = \equivrep{a,b}{}$ kanggo sawatara $a,b \in
\Nat$. Netepna $j = \equivrep{b,a}{} \in \Int$. Kanthi nggunakake tumindake
wilangan natural, $(a+b) + 0 = 0 + (a+b)$, mula $\tuple{a+b, b+a}
\sim_\Int \tuple{0,0}$ miturut definisi, lan mulane
$\equivrep{a+b, b+a}{} = \equivrep{0, 0}{} = 0_\Int$. Dadi saiki $i + j =
\equivrep{a,b}{}+\equivrep{b,a}{}=\equivrep{a+b, b+a}{}= \equivrep{0,
0}{} = 0_\Int$.
\end{proof}

Lan iki bukti \emph{Distributivitas}:

\begin{proof}
Kaya ing ndhuwur, tetepna $i = \equivrep{a_1, b_1}{}$ lan $j =
\equivrep{a_2,b_2}{}$ lan $k = \equivrep{a_3, b_3}{}$. Saiki:
\begin{align*}
	i \times (j + k)
	&= \equivrep{a_1, b_1}{} \times (\equivrep{a_2,b_2}{} + \equivrep{a_3, b_3}{})\\
	&= \equivrep{a_1, b_1}{} \times \equivrep{a_2 + a_3,b_2+b_3}{}\\
	&= \equivrep{a_1  (a_2 + a_3) + b_1  (b_2+b_3), a_1  (b_2 + b_3) + b_1 (a_2 + a_3)}{}\\
	&= \equivrep{a_1 a_2 + a_1a_3 + b_1 b_2+b_1b_3, a_1 b_2 + a_1b_3 + a_2b_1 + a_3b_1}{}\\
%		&= \equivrep{a_1 a_2 + b_1b_2 + a_1a_3 + b_1 b_2+b_1b_3, a_1 b_2 + a_1b_3 + a_2b_1 + a_3b_1}{}\\
	&= \equivrep{a_1a_2 + b_1b_2, a_1b_2 + a_2b_1}{} + \equivrep{a_1a_3 + b_1b_3, a_1b_3 + a_3b_1}{}\\
	&= (\equivrep{a_1, b_1}{} \times \equivrep{a_2,b_2}{}) + (\equivrep{a_1, b_1}{} \times  \equivrep{a_3, b_3}{})\\
	&= (i \times j) + (i \times k)
\end{align*}
\end{proof}

Kita pasrahake limang sarat liyane minangka gladhen. Sawise iku rampung,
kita wis nuduhake yen $\Int$ mbentuk ring komutatif, yaiku panambahan lan
ping-pingan kang ditetepake tumindak kaya kang kudune.

\begin{prob}
Buktekna yen $\Int$ iku ring komutatif.
\end{prob}

Nanging tugas kita durung rampung. Saliyane netepake panambahan lan
ping-pingan ing $\Int$, kita uga netepake relasi tatanan $\leq$, lan kudu
mriksa yen relasi iki tumindak kaya kang kudune. Kanthi luwih rinci, kita
kudu nuduhake yen $\Int$ mbentuk ring \emph{katata}.\footnote{Elinga saka
	\olref[sfr][rel][ord]{def:linearorder} yen tatanan total iku relasi kang
	refleksif, transitif, antisimetris, lan koneks. Ing konteks relasi tatanan,
	koneksitas kadhangkala diarani \emph{trikotomi}, amarga kanggo saben $a$
	lan $b$ kita duwe $a \leq b \lor a = b \lor a \geq b$.}

\begin{defn}
\emph{Ring katata} yaiku ring komutatif kang uga dilengkapi relasi
tatanan total $\leq$, saengga:
\begin{align*}
	a \leq b &\lif a + c \leq b + c\\
	(a \leq b \land 0 \leq c) &\lif a \times c \leq b \times c
\end{align*}
\end{defn}

\begin{prob}
Buktekna yen $\Int$ iku ring katata.
\end{prob}

Kaya sadurunge, nuduhake yen $\Int$ kang dibangun iki dadi ring katata
iku ngrepotake nanging lumrah. Pamriksane dipasrahake marang pamaca.

Iki ngrampungake prakara wilangan wutuh. Nanging saiki kita kudu nuduhake
prakara kang meh padha tumrap wilangan rasional. Mligine, kita kudu nuduhake
yen wilangan rasional mbentuk \emph{lapangan} katata miturut definisi
$+$, $\times$, lan $\leq$ kang wis diwenehake:
\begin{defn}\ollabel{orderedfield}
\emph{Lapangan katata} yaiku ring katata kang uga nyukupi:
\begin{align*}
	\emph{Invers Multiplikatif}& & (\forall a \in S \setminus \{0\})(\exists b \in S) a\times b& = 1
\end{align*}
\end{defn}

Sawise nuduhake yen $\Int$ mbentuk ring katata, nuduhake yen $\Rat$
mbentuk lapangan katata iku gampang nanging ngrepotake.

\begin{prob}
Buktekna yen $\Rat$ iku lapangan katata.
\end{prob}

Sawise ngrampungake wilangan wutuh lan rasional, mung wilangan real kang
isih kudu dirampungake. Mligine, kita kudu nuduhake yen $\Real$ mbentuk
lapangan katata kang \emph{jangkep}, yaiku lapangan katata kanthi Sipat
Kajangkepan. \olref[cuts]{realcompleteness} wis netepake yen $\Real$
nduweni Sipat Kajangkepan. Nanging kita isih kudu nindakake pamriksa kang
ngrepotake kanggo nuduhake yen $\Real$ iku lapangan katata. Cathetan
koreksi sumber OLP-0047: ukara Inggris ngilangi tembung sawisé ``the
(tedious)'`; terjemahan iki mulihake gagasan pamriksa kang dibutuhake.

Sadurunge langsung nindakake gladhen \emph{kasebut} kang ngrepotake, kita
kudu mriksa sawetara prakara kang luwih ``langsung''. Umpamane, kita butuh
jaminan yen $\alpha + \beta$ kang ditetepake iki pancen dadi
\emph{potongan} kanggo sembarang potongan $\alpha$ lan $\beta$. Iki buktine:

\begin{proof}
Amarga $\alpha$ lan $\beta$ kalorone potongan, $\alpha + \beta = \Setabs{p
+ q}{p \in \alpha \land q \in \beta}$ iku subhimpunan murni kang ora
kosong saka $\Rat$. Saiki upamakna $x < p + q$ kanggo sawatara $p \in
\alpha$ lan $q \in \beta$. Banjur $x - p < q$, mula $x - p \in \beta$,
lan $x = p + (x - p) \in \alpha + \beta$. Dadi $\alpha + \beta$ iku
segmen wiwitan saka $\Rat$. Pungkasan, kanggo sembarang $p + q \in \alpha
+ \beta$, amarga $\alpha$ lan $\beta$ kalorone potongan, ana $p_1 \in
\alpha$ lan $q_1 \in \beta$ saengga $p < p_1$ lan $q < q_1$; mula $p + q
< p_1 + q_1 \in \alpha + \beta$; dadi $\alpha + \beta$ ora duwe maksimum.
\end{proof}

Kanthi upaya kang padha, kita bisa mriksa yen $\alpha - \beta$,
$\alpha \times \beta$, lan $\alpha \div \beta$ iku potongan (kanggo kasus
pungkasan, kejaba nalika $\beta$ iku potongan nol). Nanging iki uga mung
dipasrahake marang pamaca. Cathetan sumber OLP-0047: rumus pangurangan lan
pambagian kang dirujuk ing kene mung ana minangka komentar kang dipateni
ing perangan potongan; gladhen iki nganggep definisi kang dimaksud mau.

\begin{prob}
Buktekna yen $\Real$ iku lapangan katata.
\end{prob}

Nanging isih ana prakara cilik kang kudu dirapekake. Ing
\olref[cuts]{sec}, kita ngusulake yen $\sqrt{2} =
\Setabs{p \in \Rat}{p < 0 \text{ utawa }p^2 < 2}$. Nanging kita pancen
kudu nuduhake yen himpunan iki \emph{potongan}. Iki buktine:

\begin{proof}
Cetha yen iki segmen wiwitan murni kang ora kosong saka wilangan rasional;
dadi cukup nuduhake yen ora duwe maksimum. Mligine, cukup nuduhake yen
nalika $p$ iku wilangan rasional positif kanthi $p^2 < 2$ lan $q =
\frac{2p+2}{p+2}$, kita duwe $p < q$ lan $q^2 < 2$. Kanggo ndeleng yen
$p < q$, gatekna:
\begin{align*}
	p^2 &< 2\\
	p^2 + 2p &< 2 + 2p\\
	p(p + 2) &< 2 + 2p\\
	p &< \tfrac{2+2p}{p+2} = q
\end{align*}
Kanggo ndeleng yen $q^2 < 2$, gatekna:
\begin{align*}
	p^2 &< 2\\
	2p^2 + 4p + 2 &< p^2 + 4p+ 4\\
	4p^2 + 8p + 4 &< 2(p^2 + 4p + 4)\\
	(2p+2)^2 & <2(p+2)^2\\
	\tfrac{(2p+2)^2}{(p+2)^2} &< 2\\
	q^2 &< 2
\end{align*}
\end{proof}
\end{document}
